\documentclass[11pt]{article}
\usepackage[a4paper,margin=2.5cm]{geometry}
\usepackage{amsmath,amssymb}
\usepackage{graphicx}
\usepackage{booktabs}
\usepackage{xcolor}
\usepackage{tikz}
\usepackage[colorlinks=true,linkcolor=blue!60!black,citecolor=teal]{hyperref}

\title{LaTeX Live Preview\\\large A demo document}
\author{Extension playground}
\date{\today}

\begin{document}
\maketitle

\begin{abstract}
This document exercises common packages: \texttt{amsmath}, \texttt{graphicx},
\texttt{booktabs}, \texttt{tikz}, \texttt{hyperref} and a Bib\TeX{} bibliography.
Edit any file and watch the preview update.
\end{abstract}

% !TEX root = ../main.tex
\section{Introduction}
Overleaf and the Typst web app made a tight edit--preview loop the norm.
This extension brings the same loop into the editor: the PDF on the right
is recompiled from your \emph{unsaved} buffers a moment after you stop typing,
and a double-click on the PDF jumps back to the source line.

\subsection{Why Tectonic?}
Tectonic is a self-contained \TeX{} engine. It needs no system \TeX{} Live
installation and downloads only the packages a document actually uses.


\section{Mathematics}
The Gaussian integral is therefore a very impotant thingy
\begin{equation}
  \int_{-\infty}^{\infty} e^{-x^2}\,\mathrm{d}x = \sqrt{\pi}.
  \label{eq:gauss}
\end{equation}
Equation~\eqref{eq:gauss} is classic; see also Knuth~\cite{knuth1984}.

\begin{align}
  f(x) &= \sum_{n=0}^{\infty} \frac{f^{(n)}(a)}{n!}(x-a)^n \\
  \nabla \times \mathbf{B} &= \mu_0 \mathbf{J} + \mu_0\varepsilon_0 \frac{\partial \mathbf{E}}{\partial t}
\end{align}

\section{Tables and figures}
\begin{table}[h]
  \centering
  \begin{tabular}{lrr}
    \toprule
    Engine & Cold start & Warm rebuild \\
    \midrule
    Tectonic & 60--100 s & 0.5 s \\
    latexmk + pdflatex & 2 s & 1 s \\
    \bottomrule
  \end{tabular}
  \caption{Rough compile times for a small document.}
\end{table}

\begin{figure}[h]
  \centering
  \includegraphics[width=0.35\linewidth]{figures/gradient.png}\hfill
  \begin{tikzpicture}[scale=1.1]
    \draw[->] (-0.2,0) -- (3.2,0) node[right] {$x$};
    \draw[->] (0,-0.2) -- (0,2.2) node[above] {$y$};
    \draw[thick,blue,domain=0:3,smooth] plot (\x,{0.2*\x*\x});
    \fill[red] (2,0.8) circle (2pt) node[above left] {$P$};
  \end{tikzpicture}
  \caption{A raster image (left) and a TikZ drawing (right).}
\end{figure}

% !TEX root = ../main.tex
\section{Conclusion}
Lamport's \LaTeX{} book~\cite{lamport1994} remains a good reference.
\begin{itemize}
  \item Live preview while typing.
  \item Errors in the Problems panel.
  \item One-click PDF export.
\end{itemize}


\bibliographystyle{plain}
\bibliography{refs}
\end{document}
